\documentclass[12pt,a4paper]{article}
\usepackage[utf8]{inputenc}
\usepackage[T1]{fontenc}
\usepackage[indonesian]{babel}
\usepackage{geometry}
\usepackage{longtable}
\usepackage{array}
\usepackage{booktabs}
\usepackage{enumitem}
\usepackage[hidelinks,breaklinks]{hyperref}
\geometry{margin=2.5cm}

\begin{document}

\begin{center}
\textbf{\Large RENCANA PEMBELAJARAN SEMESTER (RPS)}\\
\textbf{Berbasis Outcome-Based Education (OBE)}\\[0.5cm]
\textbf{PROGRAM STUDI TEKNIK INFORMATIKA}\\
\textbf{FAKULTAS TEKNOLOGI INFORMASI}\\
\textbf{UNIVERSITAS BALE BANDUNG}\\
\end{center}

\vspace{1cm}

% ============================================================
% 1. IDENTITAS MATA KULIAH
% ============================================================
\section*{1. Identitas Mata Kuliah}
\begin{tabular}{ll}
Nama Program Studi & : Teknik Informatika \\
Nama Mata Kuliah & : Pemrograman Web \\
Kode Mata Kuliah & : TIF-XXX \\
Semester & : 3 (Tiga) \\
SKS / Bobot Kredit & : 2 SKS \\
Jumlah Pertemuan & : 16 pertemuan \\
Komposisi Pembelajaran & : 70\% Praktik, 30\% Teori \\
Dosen Pengampu & : TBD \\
Tanggal Penyusunan & : 2 Maret 2026 \\
\end{tabular}

\vspace{0.5cm}

% ============================================================
% 2. CAPAIAN PEMBELAJARAN LULUSAN (CPL)
% ============================================================
\section*{2. Capaian Pembelajaran Lulusan (CPL)}

CPL yang dibebankan pada mata kuliah ini mencakup kompetensi lulusan dalam aspek pengetahuan, keterampilan, dan sikap:

\begin{itemize}[leftmargin=*]
  \item \textbf{CPL-04:} Kemampuan mendesain, mengimplementasi dan mengevaluasi solusi berbasis computing yang memenuhi kebutuhan-kebutuhan computing pada sebuah disiplin program.
  \item \textbf{CPL-05:} Mengidentifikasi dan menganalisis kebutuhan-kebutuhan pengguna dan mempertimbangkannya dalam memilih, membuat, mengintegrasi, mengevaluasi, dan mengadministrasi sistem berbasis computing.
\end{itemize}

\vspace{0.5cm}

% ============================================================
% 3. CAPAIAN PEMBELAJARAN MATA KULIAH (CPMK)
% ============================================================
\section*{3. Capaian Pembelajaran Mata Kuliah (CPMK)}

Kemampuan atau kompetensi spesifik yang diharapkan mahasiswa kuasai setelah menyelesaikan mata kuliah:

\begin{itemize}[leftmargin=*]
  \item \textbf{CPMK-1 (Analisis Kebutuhan Pengguna dan Perancangan Antarmuka):} Mahasiswa mampu mengidentifikasi kebutuhan fungsional dan non-fungsional pengguna untuk merancang solusi web yang intuitif melalui pembuatan wireframe dan mockup (UI/UX) yang mempertimbangkan aspek aksesibilitas dan pengalaman pengguna.
  \item \textbf{CPMK-2 (Implementasi Solusi Web Front-End):} Mahasiswa mampu mendesain dan mengimplementasikan antarmuka web yang responsif dan dinamis menggunakan teknologi standar (HTML5, CSS3, dan JavaScript) serta framework modern untuk memenuhi spesifikasi perangkat pengguna yang beragam.
  \item \textbf{CPMK-3 (Pengembangan Sistem Back-End dan Integrasi Data):} Mahasiswa mampu membangun arsitektur sisi server (back-end), mengelola basis data, dan mengintegrasikan API untuk menciptakan sistem berbasis computing yang fungsional dan aman.
  \item \textbf{CPMK-4 (Evaluasi, Deployment dan Administrasi Web):} Mahasiswa mampu melakukan pengujian, mengevaluasi keamanan serta performa aplikasi web, dan melakukan proses deployment ke server produksi untuk memastikan aplikasi dapat diakses secara optimal oleh publik.
\end{itemize}

\vspace{0.5cm}

% ============================================================
% 4. SUB-CPMK / INDIKATOR PENCAPAIAN
% ============================================================
\section*{4. Sub-CPMK / Indikator Pencapaian}

Penjabaran CPMK menjadi indikator yang lebih terukur dan dapat diuji:

\textbf{Turunan CPMK-1}
\begin{itemize}[leftmargin=*]
  \item \textbf{Sub-CPMK 1.1:} Mengidentifikasi kebutuhan fungsional dan non-fungsional pengguna untuk solusi web.
  \item \textbf{Sub-CPMK 1.2:} Membuat wireframe dan mockup antarmuka web yang mendukung alur pengguna (user flow).
  \item \textbf{Sub-CPMK 1.3:} Menerapkan prinsip aksesibilitas (a11y) dan pengalaman pengguna (UX) dalam desain antarmuka.
\end{itemize}

\textbf{Turunan CPMK-2}
\begin{itemize}[leftmargin=*]
  \item \textbf{Sub-CPMK 2.1:} Membuat halaman web dengan HTML5 (struktur, elemen semantik, form, media).
  \item \textbf{Sub-CPMK 2.2:} Menerapkan CSS3 untuk styling, layout (Flexbox/Grid), dan desain responsif.
  \item \textbf{Sub-CPMK 2.3:} Mengimplementasikan interaktivitas dengan JavaScript (DOM, event handling, validasi form).
  \item \textbf{Sub-CPMK 2.4:} Menggunakan library atau framework front-end untuk antarmuka dinamis.
\end{itemize}

\textbf{Turunan CPMK-3}
\begin{itemize}[leftmargin=*]
  \item \textbf{Sub-CPMK 3.1:} Mengimplementasikan logika server-side dengan PHP (pemrosesan form, struktur dasar, OOP).
  \item \textbf{Sub-CPMK 3.2:} Mengelola basis data dan mengimplementasikan operasi CRUD dengan MySQL serta koneksi dari PHP.
  \item \textbf{Sub-CPMK 3.3:} Mengintegrasikan API dan menghubungkan front-end dengan back-end (request/response, data terstruktur).
\end{itemize}

\textbf{Turunan CPMK-4}
\begin{itemize}[leftmargin=*]
  \item \textbf{Sub-CPMK 4.1:} Melakukan pengujian dan evaluasi performa aplikasi web.
  \item \textbf{Sub-CPMK 4.2:} Menerapkan prinsip keamanan web: validasi input, prepared statement, session, mitigasi XSS dan CSRF.
  \item \textbf{Sub-CPMK 4.3:} Melakukan deployment ke lingkungan produksi dan administrasi dasar (server/web, basis data).
\end{itemize}

\vspace{0.5cm}

% ============================================================
% 5. MATERI PEMBELAJARAN (BAHAN KAJIAN)
% ============================================================
\section*{5. Materi Pembelajaran (Bahan Kajian)}

Materi disusun dari sumber terbuka dan judul materi yang relevan, lalu dipetakan ke Sub-CPMK dalam 16 pertemuan.

\textbf{Rencana Pembelajaran 16 Pertemuan}

\begin{longtable}{|c|>{\raggedright\arraybackslash}p{3cm}|>{\raggedright\arraybackslash}p{6.2cm}|>{\raggedright\arraybackslash}p{2.2cm}|>{\raggedright\arraybackslash}p{2.2cm}|}
\hline
\textbf{Pertemuan} & \textbf{Sub-CPMK} & \textbf{Materi Pokok} & \textbf{Bentuk} & \textbf{Asesmen} \\
\hline
\endfirsthead
\hline
\textbf{Pertemuan} & \textbf{Sub-CPMK} & \textbf{Materi Pokok} & \textbf{Bentuk} & \textbf{Asesmen} \\
\hline
\endhead
\hline
\endfoot

1 & Sub-CPMK 1.1 & Pengantar pemrograman web, arsitektur client-server, ruang lingkup kebutuhan pengguna, kontrak kuliah & Teori & Kuis awal \\
\hline
2 & Sub-CPMK 1.1, 1.2 & Kebutuhan fungsional dan non-fungsional, user flow, wireframe dasar & Praktik & Tugas 1 \\
\hline
3 & Sub-CPMK 1.2, 1.3 & Mockup, prototyping, aksesibilitas dasar (WCAG), prinsip UX & Praktik & Tugas 2 \\
\hline
4 & Sub-CPMK 2.1 & HTML dasar, struktur dokumen, elemen semantik & Praktik & Praktikum 1 \\
\hline
5 & Sub-CPMK 2.1, 1.3 & HTML form, media, validasi dasar, aksesibilitas form & Praktik & Praktikum 2 \\
\hline
6 & Sub-CPMK 2.2 & CSS dasar, selector, box model, tipografi & Praktik & Praktikum 3 \\
\hline
7 & Sub-CPMK 2.2 & Flexbox, Grid, desain responsif & Praktik & Tugas 3 \\
\hline
8 & Sub-CPMK 1.1-2.2 & UTS (teori dan praktik) & Evaluasi & UTS \\
\hline
9 & Sub-CPMK 2.3 & JavaScript dasar, DOM, manipulasi elemen & Praktik & Praktikum 4 \\
\hline
10 & Sub-CPMK 2.3 & Event handling, validasi form, error handling & Praktik & Praktikum 5 \\
\hline
11 & Sub-CPMK 2.4 & Library/framework front-end (jQuery/Bootstrap/Vue/React) untuk UI dinamis & Praktik & Tugas 4 \\
\hline
12 & Sub-CPMK 3.1 & PHP dasar, pemrosesan form, include/require & Praktik & Praktikum 6 \\
\hline
13 & Sub-CPMK 3.1, 4.2 & PHP OOP dasar, session/cookie, validasi input & Praktik & Praktikum 7 \\
\hline
14 & Sub-CPMK 3.2 & MySQL dasar, CRUD, koneksi PHP-MySQL (MySQLi/PDO) & Praktik & Praktikum 8 \\
\hline
15 & Sub-CPMK 3.3, 4.1, 4.2 & Integrasi API, pengujian, performa, keamanan (XSS/CSRF) & Praktik & Proyek \\
\hline
16 & Sub-CPMK 4.3 & Deployment dasar, administrasi web, presentasi hasil & Evaluasi/Praktik & UAS \\
\hline
\end{longtable}

\vspace{0.5cm}

% ============================================================
% 6. METODE PEMBELAJARAN
% ============================================================
\section*{6. Metode Pembelajaran}

Strategi atau pendekatan pembelajaran dipilih sesuai OBE dengan komposisi 70\% praktik dan 30\% teori:

\begin{itemize}[leftmargin=*]
  \item \textbf{Ceramah interaktif (teori 30\%):} Penjelasan konsep inti, diskusi terarah, dan tanya jawab.
  \item \textbf{Praktikum terbimbing (praktik 70\%):} Latihan coding HTML, CSS, JavaScript, PHP, dan MySQL.
  \item \textbf{Project-based learning:} Pengembangan proyek web terintegrasi dari front-end sampai back-end.
  \item \textbf{Problem-based learning:} Analisis kebutuhan dan perancangan antarmuka dari studi kasus.
  \item \textbf{Demonstrasi dan code review:} Umpan balik cepat atas hasil praktik mahasiswa.
\end{itemize}

\vspace{0.5cm}

% ============================================================
% 7. PENGALAMAN BELAJAR MAHASISWA
% ============================================================
\section*{7. Pengalaman Belajar Mahasiswa}

Deskripsi tugas, aktivitas, atau pengalaman belajar yang mendukung pencapaian Sub-CPMK:

\begin{itemize}[leftmargin=*]
  \item Mengidentifikasi kebutuhan fungsional dan non-fungsional dari studi kasus sederhana.
  \item Menyusun user flow serta membuat wireframe dan mockup antarmuka web.
  \item Membangun halaman web semantik dan aksesibel menggunakan HTML5.
  \item Mendesain tampilan responsif dengan CSS (Flexbox/Grid).
  \item Membuat interaksi front-end dengan JavaScript (DOM dan event).
  \item Mengembangkan aplikasi web sederhana dengan PHP (form, session, dan validasi).
  \item Mengelola database MySQL dan menerapkan operasi CRUD dari PHP.
  \item Menerapkan keamanan dasar (prepared statement, XSS, CSRF).
  \item Melakukan pengujian dasar dan evaluasi performa aplikasi web.
  \item Melakukan deployment sederhana dan mempresentasikan hasil proyek.
\end{itemize}

\vspace{0.5cm}

% ============================================================
% 8. KRITERIA, INDIKATOR, DAN BOBOT PENILAIAN
% ============================================================
\section*{8. Kriteria, Indikator, dan Bobot Penilaian}

Teknik/alat asesmen dipetakan ke Sub-CPMK/CPMK dengan bobot yang jelas:

\begin{longtable}{|>{\raggedright\arraybackslash}p{2.8cm}|>{\raggedright\arraybackslash}p{4cm}|>{\raggedright\arraybackslash}p{5.7cm}|c|}
\hline
\textbf{Komponen} & \textbf{Teknik Asesmen} & \textbf{Indikator/CPMK} & \textbf{Bobot (\%)} \\
\hline
\endfirsthead
\hline
\textbf{Komponen} & \textbf{Teknik Asesmen} & \textbf{Indikator/CPMK} & \textbf{Bobot (\%)} \\
\hline
\endhead
\hline
\endfoot

Tugas/Praktikum & Lab exercise dan laporan praktik & Sub-CPMK 1.1-3.2 & 35 \\
\hline
Kuis & Kuis singkat dan refleksi & Sub-CPMK 1.1-2.2 & 10 \\
\hline
UTS & Ujian tengah semester (teori dan praktik) & CPMK-1, CPMK-2 & 20 \\
\hline
Proyek & Proyek terintegrasi front-end dan back-end & CPMK-2, CPMK-3, CPMK-4 & 20 \\
\hline
UAS & Ujian akhir dan presentasi hasil & CPMK-1, CPMK-2, CPMK-3, CPMK-4 & 15 \\
\hline
\textbf{Total} & & & \textbf{100} \\
\hline
\end{longtable}

\textbf{Kriteria Penilaian:}
\begin{itemize}[leftmargin=*]
  \item A (85-100): Menguasai semua CPMK dengan sangat baik dan mampu menerapkan dalam proyek nyata.
  \item B (70-84): Menguasai sebagian besar CPMK dengan baik.
  \item C (60-69): Menguasai CPMK dasar dengan cukup.
  \item D (50-59): Menguasai sebagian kecil CPMK.
  \item E (<50): Belum menguasai CPMK yang ditetapkan.
\end{itemize}

\vspace{0.5cm}

% ============================================================
% 9. EVALUASI DAN REFLEKSI PEMBELAJARAN (OPSIONAL)
% ============================================================
\section*{9. Evaluasi dan Refleksi Pembelajaran}

Penilaian sumatif/formatif untuk memantau ketercapaian outcome secara menyeluruh:

\begin{itemize}[leftmargin=*]
  \item \textbf{Evaluasi formatif:} Kuis mingguan, praktik, dan umpan balik cepat.
  \item \textbf{Evaluasi sumatif:} UTS dan UAS untuk mengukur capaian CPMK secara komprehensif.
  \item \textbf{Refleksi mahasiswa:} Jurnal atau ringkasan belajar untuk menilai pemahaman.
  \item \textbf{Evaluasi dosen:} Review hasil proyek dan perbaikan strategi pembelajaran.
  \item \textbf{Continuous improvement:} Analisis capaian untuk perbaikan RPS semester berikutnya.
\end{itemize}

\vspace{0.5cm}

% ============================================================
% 10. DAFTAR REFERENSI
% ============================================================
\section*{10. Daftar Referensi}

Sumber belajar utama yang digunakan dalam penyusunan materi dan asesmen:

\begin{enumerate}[leftmargin=*]
  \item \textit{RPS Pemrograman Web - Universitas Mulawarman}. \url{https://si.ft.unmul.ac.id/kurikulum_rps/1753877341_28.%20RPS-Pemrograman%20Web.pdf}
  \item \textit{RPS Pemrograman Web - UIN Ar-Raniry}. \url{https://cdn.ar-raniry.ac.id/aps-test/teknologi_informasi/Pemograman_Web.pdf}
  \item \textit{Panduan Kurikulum Berbasis OBE/KKNI/SKKNI - APTIKOM}. \url{https://aptikom-journal.id/home/index.php/2024/05/31/panduan-kurikulum-berbasis-obe-kkni-skkni-aptikom-program-studi-sarjana-teknologi-infomasi-2/}
  \item \textit{MDN Learn Web Development}. \url{https://developer.mozilla.org/en-US/docs/Learn}
  \item \textit{W3Schools Tutorials}. \url{https://www.w3schools.com/tutorials/}
  \item \textit{web.dev Learn Web Development}. \url{https://web.dev/learn}
  \item \textit{MDN - Introduction to HTML}. \url{https://developer.mozilla.org/en-US/docs/Learn/HTML/Introduction_to_HTML/Getting_started}
  \item \textit{web.dev - Learn HTML}. \url{https://web.dev/learn/html/welcome}
  \item \textit{MDN - Styling the Web with CSS}. \url{https://developer.mozilla.org/en-US/docs/Learn/CSS}
  \item \textit{web.dev - Learn CSS}. \url{https://web.dev/learn/css}
  \item \textit{JavaScript.info (Indonesia)}. \url{https://id.javascript.info/}
  \item \textit{MDN - JavaScript Guide}. \url{https://developer.mozilla.org/en-US/docs/Web/JavaScript/Guide}
  \item \textit{PHP Manual}. \url{https://www.php.net/manual/en/}
  \item \textit{PHP The Right Way}. \url{https://phptherightway.com}
  \item \textit{MySQL Documentation}. \url{https://dev.mysql.com/doc/}
  \item \textit{MySQL Getting Started Guide}. \url{https://dev.mysql.com/doc/mysql-getting-started/en/}
  \item \textit{MDN - Tools and testing}. \url{https://developer.mozilla.org/en-US/docs/Learn/Tools_and_testing}
  \item \textit{MDN - Deploying to production}. \url{https://developer.mozilla.org/en-US/docs/Learn/Server-side/Express_Nodejs/deployment}
  \item \textit{OWASP - Cross Site Scripting Prevention Cheat Sheet}. \url{https://cheatsheetseries.owasp.org/cheatsheets/Cross_Site_Scripting_Prevention_Cheat_Sheet.html}
  \item \textit{OWASP - Cross-Site Request Forgery Prevention Cheat Sheet}. \url{https://cheatsheetseries.owasp.org/cheatsheets/Cross-Site_Request_Forgery_Prevention_Cheat_Sheet.html}
\end{enumerate}

\end{document}
